\chapter{CONCLUSION AND FUTURE WORK}
\label{ch:conclusion}

In conclusion, Phase I of the project has produced a complete, tested prediction pipeline for explainable visual commonsense reasoning about social interactions. The system accepts image, video and live-camera input, standardises it with OpenCV, performs visual understanding and language reasoning with SmolVLM, enriches the detected concepts with ConceptNet, and fuses the resulting evidence in a configurable social-interaction reasoning module.

SmolVLM was fine-tuned on the VCR dataset with LoRA for 20 epochs, and the adapter from epoch 11 was selected on a held-out selection set. On separate held-out test questions it improves accuracy over the base model from 44.7\% to 58.7\% for \QtoA, from 41.3\% to 57.7\% for \QAtoR{} and from 15.3\% to 32.7\% for \QtoAR.

The main output of Phase I is a structured prediction of the activity, relationship and intention, accompanied by confidence values, alternatives, evidence and a textual explanation. The modules are connected through a reusable \code{Phase1Pipeline}, the FastAPI--React interface demonstrates the complete flow, and all 23 automated tests pass. Phase I objectives 1--10 of Section~\ref{sec:objectives} have therefore been met, and the implementation provides a tested foundation for the second phase.

As future work, Phase II will extend the system from prediction to verified explanation. Grad-CAM maps over the SmolVLM vision encoder will form a visual explanation path, and the concepts cited in the generated explanation will form an independent textual path. The Jaccard-based cross-modal consistency between the two concept sets (Equation~\ref{eq:cmc}) will be measured, faithfulness will be evaluated with the Insertion--Deletion metric and the Pointing Game, and the input frame, the predicted interaction, the three labels, the explanations and the consistency score will be presented in an explainability dashboard.

In addition, the following improvements are planned:
\begin{itemize}
\item a labelled evaluation set for the activity, relationship and intention predictions of the full pipeline, together with an ablation study of the fusion weights and of each evidence source;
\item measurement and improvement of confidence calibration;
\item a larger VCR test set and training on more VCR questions;
\item faster inference by batching the VLM prompts, to approach real-time video analysis;
\item privacy safeguards that prevent the naming of individuals, and an audit of the predictions across demographic groups.
\end{itemize}
